% Generated draft: requires mathematical and source review.
\documentclass[11pt]{article}
\usepackage[margin=25mm]{geometry}
\usepackage{fontspec}
\setmainfont{Latin Modern Roman}
\usepackage{amsmath,amssymb,mathtools,unicode-math}
\setmathfont{Latin Modern Math}
\usepackage{enumitem,needspace,xurl,hyperref}
\setlength{\parindent}{0pt}
\setlength{\parskip}{5pt}
\setlength{\emergencystretch}{3em}
\newcommand{\N}{\mathbb{N}}
\newcommand{\Z}{\mathbb{Z}}
\newcommand{\Q}{\mathbb{Q}}
\newcommand{\R}{\mathbb{R}}
\newcommand{\C}{\mathbb{C}}
\newcommand{\F}{\mathbb{F}}
\newcommand{\A}{\mathbb{A}}
\newcommand{\PP}{\mathbb{P}}
\newcommand{\E}{\mathbb{E}}
\newcommand{\eps}{\varepsilon}
\newcommand{\dd}{\,\mathrm{d}}
\begin{document}
\section*{UnsolvedMath ID: 30006566}
\textbf{Source problem code:} catalogue entry
\subsection{Definitions and mathematical statement}
The supplied record asks for two reverse implications between approximation properties of traces. It does not specify the ambient class of algebras, separability, or unitality. The following is a precise conditional formulation for separable unital complex \(C^*\)-algebras; this scope is a drafting convention requiring source verification, not an assumption recovered from the supplied excerpt.

Let \(A\) be a separable unital complex \(C^*\)-algebra. A tracial state on \(A\) is a positive linear functional \(\tau:A\to\C\) satisfying \(\tau(1_A)=1\) and \(\tau(ab)=\tau(ba)\) for all \(a,b\in A\). For \(x\in M_n(\C)\), put
\[
\operatorname{tr}_n(x)=\frac{1}{n}\operatorname{Tr}(x),\qquad
\|x\|_{2,n}=\operatorname{tr}_n(x^*x)^{1/2}.
\]
Choose positive integers \(n_k\) and a nonprincipal ultrafilter \(\omega\) on \(\N\). Let
\[
B=\left\{(x_k)\in\prod_{k\in\N}M_{n_k}(\C):\sup_k\|x_k\|<\infty\right\},
\]
\[
I_{\|\cdot\|,\omega}=\{(x_k)\in B:\lim_{k\to\omega}\|x_k\|=0\},\qquad
I_{2,\omega}=\{(x_k)\in B:\lim_{k\to\omega}\|x_k\|_{2,n_k}=0\}.
\]
Each quotient \(B/I_{\|\cdot\|,\omega}\) and \(B/I_{2,\omega}\) carries the trace
\[
\operatorname{tr}_{\omega}([(x_k)])=\lim_{k\to\omega}\operatorname{tr}_{n_k}(x_k).
\]
Under the conventions used here, \(\tau\) is an \emph{MF trace} if, for some such choices, there is a unital \(*\)-homomorphism
\[
\pi:A\longrightarrow B/I_{\|\cdot\|,\omega},\qquad
\tau=\operatorname{tr}_{\omega}\circ\pi.
\]
It is \emph{hyperlinear} if the corresponding condition holds with target \(B/I_{2,\omega}\). These homomorphisms are not required to be injective; faithfulness of \(\tau\) is not assumed.

The trace \(\tau\) is \emph{amenable} if, for every finite subset \(F\subseteq A\) and every \(\varepsilon>0\), there exist \(n\geq1\) and a unital completely positive linear map \(\phi:A\to M_n(\C)\) such that
\[
|\operatorname{tr}_n(\phi(a))-\tau(a)|<\varepsilon\quad(a\in F),
\]
\[
\|\phi(ab)-\phi(a)\phi(b)\|_{2,n}<\varepsilon\quad(a,b\in F).
\]
Here complete positivity means that \(\operatorname{id}_{M_m}\otimes\phi\) is positive for every \(m\geq1\). The trace is \emph{quasidiagonal} if the same finite-set condition holds with the operator-norm estimate
\[
\|\phi(ab)-\phi(a)\phi(b)\|<\varepsilon\quad(a,b\in F)
\]
in place of the normalized \(2\)-norm estimate.

With these conventions, the two questions are whether, for every separable unital complex \(C^*\)-algebra \(A\) and every tracial state \(\tau\) on \(A\),
\[
\tau\text{ hyperlinear}\ \Longrightarrow\ \tau\text{ an MF trace},
\]
and, separately,
\[
\tau\text{ amenable}\ \Longrightarrow\ \tau\text{ quasidiagonal}.
\]
The supplied curation describes the forward implications from MF to hyperlinear and from quasidiagonal to amenable, and labels the reverse implications open. The intended algebra class and the report's precise terminology must be checked before this conditional formulation can be treated as the source's fully quantified statement.
\subsection{Short English statement}
The questions ask whether every hyperlinear trace is MF and whether every amenable trace is quasidiagonal. A conditional formulation for separable unital complex \(C^*\)-algebras is given above; the supplied excerpt does not establish that scope. The open status is supplied dated curation, not an independently verified assessment.
\subsection{Sources}
\begin{itemize}
\item Local UnsolvedMath registry, permanent ID \(30006566\), title \emph{Reverse Approximation Implications for Traces}, source problem code \texttt{catalogue entry}. The supplied record belongs to the Oberwolfach Reports Open Problems collection and gives proposed year 2026; no individual proposer is supplied.
\item \emph{Noncommutative Harmonic Analysis and Quantum Information} (2026), Oberwolfach Reports, pp.\ 217--275, DOI \texttt{10.4171/owr/2026/4}. Supplied source and literature reference: \url{https://doi.org/10.4171/owr/2026/4}, also labeled ``Oberwolfach Report 4/2026.'' The original extracted fragment says that the opposite implications are open and includes the report title and page number 227, but omits the preceding implications.
\item The registry marks the statement \texttt{corrected\_verified} and reports that it was reconstructed and checked against the official report. This is supplied provenance; no independent source visit or verification is asserted here.
\item Literature assessment dated exactly \texttt{2026-08-22}: the supplied curation states that MF traces imply hyperlinearity and quasidiagonal traces imply amenability, while the reverse implications remain open. This is supplied dated curation derived from the OWR record, not independently verified literature status. The supplied background explicitly calls for independent verification of time-sensitive status and cited claims.
\end{itemize}
\end{document}
